\section{Experimental Evaluation}
\label{sec:exp}

We evaluate our method with three research questions.
\textbf{RQ1 (Effectiveness)} asks whether it makes the parser more
accurate than existing clarification methods.
\textbf{RQ2 (Fewer questions)} asks whether it can ask fewer questions
without losing accuracy. \textbf{RQ3 (Real users)} asks whether the
gains hold when real users answer the questions.

%======================================================================
\subsection{Experimental Setup}
\label{sec:exp-setup}

\textbf{Benchmarks and metrics.} \emph{BIRD}~\cite{dataset-bird}: the 498 questions
of BIRD Mini-Dev with the corrected annotations of
Arcwise~\cite{arcwise}, over real databases and with an official
evidence for each question. \emph{AMBROSIA}~\cite{ambrosia}: the
1{,}149 ambiguous test questions, involving scope, attachment, and
vagueness, with one gold query per interpretation (two for 779
questions, three for 370). \emph{CLAMBSQL}~\cite{zhang2025clear}:
200 questions for each of seven ambiguity types drawn from AmbiQT
and AMBROSIA, with two or three gold queries each. As the intended
query $y^{*}$, we take the first AMBROSIA gold query in the released
order and the gold query designated by CLAMBSQL. AMBROSIA and
CLAMBSQL have no evidence, so on them all evidence that the parser
receives comes from the clarifications. We report execution accuracy
(EX), where a predicted query is correct if its result on $D$ equals
that of the gold query, compared as sets with their column order, and
the average number of questions per instance ($Q$). For our methods,
$Q$ counts ambiguities sent to the user, and automatic resolutions are
not counted. For CLEAR and AmbiSQL, $Q$ counts the clarification slots
each method asks. We test differences in EX with the exact McNemar
test (two-sided).

The rest of the setup follows the pipeline of
Fig.~\ref{fig:overview}. Our method detects, checks, and scores the
ambiguities, a simulated user answers the questions that remain, and
a parser writes the final query from the clarified input.

\textbf{Clarification models.} Our method uses two models.
Detection, the ambiguity checks, and the relevance check run on
\coder, which denotes Qwen3-Coder-30B-A3B-Instruct~\cite{qwen3}
served locally, whichever parser writes the final query. The relevance
check therefore asks whether an answer changes the query that \coder{}
writes, and for the other parsers it uses this as an approximation of
their behavior. The scorer is Qwen3-8B, fine-tuned with
LoRA~\cite{lora} (rank 16, scaling factor 32) on \trainCards{}
ambiguities that our pipeline produces on \trainCardQuestions{}
questions from 60 BIRD training databases, labeled by the simulated
user. \valCards{} ambiguities from \valCardQuestions{} questions on
nine further training databases serve as validation data. The split
is by database, and none of these databases appears in the three
benchmarks. We select $\theta$ on the validation ambiguities by the
agreement of Eq.~\eqref{eq:theta} with $\alpha=0.9$, over
$\theta\in\{0.1,0.2,\dots,0.9\}$. The selected threshold is
$\theta^{*}=\thetaJev$, at which the scorer auto-resolves 305
validation ambiguities (\abValAuto) and agrees with the simulated user
on \abValAgree{} of them. The same $\theta^{*}$ is applied unchanged
to all benchmarks and parsers. Ambiguities with more than 26
interpretations are treated as $H_{\mathrm{norm}}=1$.

\textbf{Simulated user.} As in prior
work~\cite{zhang2025clear,xu2026visinteract}, the simulated user knows
the intended query. It is \coder{} given the gold query and the
ambiguity, and it selects the interpretation consistent with the gold
query, or answers ``Others'' if no interpretation is. Two
checks support this proxy. On the \simN{} CLAMBSQL spans whose
annotated gold column is among our interpretations, it selects the
gold choice on \simPct{} (\simHit/\simN). With \terra{} as the
simulated user instead, no EX changes significantly (smallest
$p=\pSimTeClambsqlAskAll$). Section~\ref{sec:exp-user} replaces it
with real users.

\textbf{Parsers.} Our method only appends clarifications to the
evidence and can be combined with any parser. We use three parsers of
increasing strength. \coder{} is prompted zero-shot with the question,
the evidence, and the table and column names. DeepEye-SQL~\cite{deepeye}
runs with its official pipeline and \coder{} as its backbone.
\terra{} is a closed-source model of OpenAI's GPT-5.6
family~\cite{openai2026gpt56}, called with the zero-shot prompt of
\coder, \texttt{reasoning\_effort=\terraEffort}, and the default
temperature on \terraCallDate. BIRD may be part of its training data,
which affects all systems equally because they share the parser.

\textbf{Baselines.} \emph{Never-Ask} runs the parser on the original
input. \emph{CLEAR}~\cite{zhang2025clear} and
\emph{AmbiSQL}~\cite{ding2025ambisql} detect ambiguities before
parsing and ask the user about every detected one. We re-implement
both from their papers and code, keeping their detection prompts and
post-detection steps. CLEAR by design rewrites the question and the
schema instead of only adding evidence, and we keep this behavior.
AmbiSQL answers keep their native form, which may select several
mappings (\ambiMultiPctMin--\ambiMultiPctMax{} of questions). All
systems share the detection LLM, the simulated user, and the parser,
so differences in EX come from what is detected and how it is
resolved. \emph{Ours} is the full pipeline, in which the scorer
decides which of the ambiguities that survive the checks are written
into the evidence and which are asked. \emph{Ours w/o scorer} asks
about every ambiguity that survives the checks.

\textbf{Implementation details.} Every system is run once, and all
variants of our method share one detection output per benchmark, as
well as the simulated responses and parser outputs, which are cached
by the final evidence string. Whenever the parser input of any system
(question, schema, and evidence) is identical to the
no-clarification input, we reuse the no-clarification prediction, so
differences between systems are caused only by the clarifications they
write. All open-weight models use greedy decoding with thinking
disabled, output is limited to 4{,}096 tokens for our pipeline, and
queries time out after 30 seconds.

%======================================================================
\subsection{Overall Effectiveness (RQ1)}
\label{sec:exp-rq1}

\begin{table}[t]
\centering
\caption{EX (\%) and questions per instance ($Q$) with \coder{} on BIRD~\cite{dataset-bird}, AMBROSIA~\cite{ambrosia}, and CLAMBSQL~\cite{zhang2025clear}. Arrows: change over Never-Ask. Bold: best EX. Differences are computed before rounding.}
\label{tab:main}
\footnotesize
\setlength{\tabcolsep}{2.5pt}
\begin{tabular}{@{}lcccccc@{}}
\toprule
& \multicolumn{2}{c}{BIRD} & \multicolumn{2}{c}{AMBROSIA} &
\multicolumn{2}{c}{CLAMBSQL} \\
\cmidrule(lr){2-3}\cmidrule(lr){4-5}\cmidrule(lr){6-7}
Method & EX & $Q$ & EX & $Q$ & EX & $Q$ \\
\midrule
Never-Ask & \abBirdZero & 0 & \abAmbrosiaZero & 0 & \abClambsqlZero & 0 \\
CLEAR~\cite{zhang2025clear} & \abBirdClear & \qTabBirdClear & \abAmbrosiaClear & \qTabAmbrosiaClear & \abClambsqlClear & \qTabClambsqlClear \\
AmbiSQL~\cite{ding2025ambisql} & \abBirdAmbiSql & \qTabBirdAmbiSql & \abAmbrosiaAmbiSql & \qTabAmbrosiaAmbiSql & \abClambsqlAmbiSql & \qTabClambsqlAmbiSql \\
\midrule
Ours w/o scorer & \textbf{\abBirdAskAll} & \abQBirdAskAll & \abAmbrosiaAskAll & \abQAmbrosiaAskAll & \textbf{\abClambsqlAskAll} & \abQClambsqlAskAll \\
Ours & \textbf{\abBirdJevAsk} & \abQBirdJevAsk & \textbf{\abAmbrosiaJevAsk} & \abQAmbrosiaJevAsk & \abClambsqlJevAsk & \abQClambsqlJevAsk \\
\bottomrule
\end{tabular}
\end{table}

\textbf{Accuracy and questions.} Table~\ref{tab:main} compares all
systems with \coder{} as the parser. Asking every ambiguity that survives
the checks (Ours w/o scorer) improves EX over Never-Ask by 7.4, 34.6,
and 11.4 points on BIRD, AMBROSIA, and CLAMBSQL. CLEAR improves EX by
0.8, 4.8, and 2.2 points, and AmbiSQL by 0.6, 8.0, and 3.1. The
scorer auto-resolves part of the ambiguities, so Ours asks 35\%, 29\%,
and 32\% fewer questions than Ours w/o scorer, and its EX differs by
$0.0$, $+0.8$, and $-0.4$ points, none of them significant
($p=1.0$, $0.16$, and $0.33$, exact McNemar). Ours asks about as many
questions as AmbiSQL (\abQBirdJevAsk, \abQAmbrosiaJevAsk, and
\abQClambsqlJevAsk{} against \qTabBirdAmbiSql, \qTabAmbrosiaAmbiSql,
and \qTabClambsqlAmbiSql) and is 6.8, 27.4, and 7.9 points more
accurate ($p=0.002$, $3.4\times10^{-53}$, and $1.5\times10^{-11}$). The gain over AmbiSQL therefore comes
from which ambiguities are detected and how the answers are written
into the evidence.

\textbf{Detection coverage.} To test whether the gain comes from what
is detected, we count on CLAMBSQL how many of its annotated ambiguous
spans each method covers (Table~\ref{tab:recall}). CLAMBSQL annotates 1{,}690 spans on its
1{,}400 questions, since a question may carry several. A gold span is
covered if a span of the method matches it by word-boundary substring,
content-word overlap, or a fuzzy ratio of at least 0.8. Matching
ignores the type, so the three taxonomies need no alignment. Predicted
spans that cover at least 70\% of the content words of the question
are excluded, and for AmbiSQL, which outputs no spans, we take the
phrase that each of its questions quotes. A manual check of 90
automatic labels agrees with all of them. For Ours,
we take the ambiguities that survive the checks, i.e., those of Ours
w/o scorer. The scorer only decides whether each of them is asked or
resolved automatically, so the full method handles the same set.
Ours covers 72.8\% of the spans,
against 59.5\% for CLEAR and 39.7\% for AmbiSQL. It leads on five of
the seven types, most clearly on Scope (97.5\% against 50.0\% and
50.5\%).
CLEAR leads on Join and Aggregate. CLAMBSQL's mappings were initialized
with CLEAR's own detection before manual revision, so its spans follow
CLEAR's granularity, e.g., one span per aggregate word in ``average,
minimum, and maximum age''. The checks cost little coverage: before
them, Ours covers 75.1\%. Covering more costs few additional
questions: Ours asks \abQClambsqlJevAsk{} questions per instance,
against \qTabClambsqlAmbiSql{} for AmbiSQL (Table~\ref{tab:main}).

\begin{table}[t]
\centering
\caption{Span-level detection recall on CLAMBSQL~\cite{zhang2025clear} (\%). $n$: gold spans. Ours: ambiguities after the checks. \#Spans: distinct predicted spans per question, not comparable to $Q$.}
\label{tab:recall}
\footnotesize
\setlength{\tabcolsep}{4pt}
\begin{tabular}{@{}lrccc@{}}
\toprule
Gold type & $n$ & CLEAR~\cite{zhang2025clear} & AmbiSQL~\cite{ding2025ambisql} & Ours \\
\midrule
Scope       & 200 & 50.0 & 50.5 & \textbf{97.5} \\
Attachment  & 200 & 48.5 & 77.5 & \textbf{83.5} \\
Column      & 200 & 72.5 & 32.0 & \textbf{83.0} \\
Table       & 313 & 56.5 & 34.2 & \textbf{67.4} \\
Join        & 200 & \textbf{75.0} & 23.5 & 71.5 \\
Aggregate   & 377 & \textbf{66.8} & 28.6 & 59.4 \\
Vagueness   & 200 & 42.0 & 44.5 & \textbf{62.5} \\
\midrule
Overall     & 1{,}690 & 59.5 & 39.7 & \textbf{72.8} \\
\#Spans     & --  & 1.15 & 1.12 & 1.87 \\
\bottomrule
\end{tabular}
\end{table}
\begin{table}[t]
\centering
\caption{EX (\%) with strong parsers, using the ambiguities, answers, and $\theta^{*}$ of \coder{}. Arrows: change over Never-Ask. Bold: best EX per parser.}
\label{tab:terra}
\footnotesize
\begin{tabular}{@{}lcc@{}}
\toprule
Method & AMBROSIA & CLAMBSQL \\
\midrule
\multicolumn{3}{@{}l}{\emph{\terra}~\cite{openai2026gpt56}} \\
Never-Ask & 25.8 & 52.9 \\
CLEAR & 28.9\dUp{3.0} & 54.1\dUp{1.2} \\
AmbiSQL & 33.3\dUp{7.5} & 56.1\dUp{3.2} \\
Ours w/o scorer & \textbf{58.6}\dUp{32.8} & \textbf{61.4}\dUp{8.5} \\
Ours & 58.4\dUp{32.6} & 61.3\dUp{8.4} \\
\midrule
\multicolumn{3}{@{}l}{\emph{DeepEye-SQL}~\cite{deepeye}} \\
Never-Ask & 27.2 & 54.1 \\
CLEAR & 32.8\dUp{5.6} & 56.4\dUp{2.3} \\
AmbiSQL & 36.3\dUp{9.1} & 58.9\dUp{4.9} \\
Ours w/o scorer & 47.6\dUp{20.4} & \textbf{61.4}\dUp{7.3} \\
Ours & \textbf{47.7}\dUp{20.5} & 60.8\dUp{6.7} \\
\bottomrule
\end{tabular}
\end{table}
\textbf{Strong parsers.} Table~\ref{tab:terra} repeats the comparison
with \terra{} and DeepEye-SQL as the parser. With both, Ours w/o
scorer gains the most over Never-Ask, so a stronger parser does not
make the clarifications redundant. On AMBROSIA and CLAMBSQL, Ours is 25.1
and 5.1 points more accurate than AmbiSQL with \terra{}
($p=2.5\times10^{-38}$ and $2.9\times10^{-6}$) and 11.4 and 1.9
points with DeepEye-SQL ($p=4.5\times10^{-13}$ and $0.062$), so only
the last gain is not significant. The
scorer, whose $\theta^{*}$ is not tuned per parser, changes EX by at
most 0.6 points against Ours w/o scorer, never significantly
($p\ge0.077$).

%======================================================================
\subsection{Fewer Questions (RQ2)}
\label{sec:exp-rq2}

We answer RQ2 in three steps: how accuracy changes as the scorer
resolves more ambiguities, which types the scorer resolves, and which
components remove questions.

\textbf{Operating points.} Figure~\ref{fig:pareto} traces the scorer over
$\theta\in\{0.1,\dots,0.9\}$ against the baselines.

\begin{figure*}[t]
\centering
\includegraphics[width=\textwidth]{figures/pareto.pdf}
\caption{Accuracy versus questions (\coder). The curve varies $\theta$ from 0.1, where nothing is auto-resolved (Ours w/o scorer), to 0.9. Star: $\theta^{*}$.}
\label{fig:pareto}
\end{figure*}

Every operating point is more accurate than CLEAR and AmbiSQL on all
three benchmarks. At $\theta=0.9$, Ours asks 0.64, 0.52, and 0.55
questions per instance on BIRD, AMBROSIA, and CLAMBSQL, fewer than
AmbiSQL everywhere and fewer than CLEAR on AMBROSIA and CLAMBSQL, and
is still 5.2, 23.6, and 5.4 points more accurate than the stronger of
the two. Up to $\theta^{*}$, EX stays within 0.8 points of Ours w/o
scorer. Beyond $\theta^{*}$, it falls by up to 1.6, 3.8, and 2.5 points, so
$\theta^{*}$ sits where $Q$ has dropped most before accuracy starts to
fall. At $\theta^{*}$, the auto-resolved interpretations agree with
the simulated user on 92\%, 94\%, and 67\% of the ambiguities, against
\abValAgree{} on validation. On CLAMBSQL, whose readings are
constructed to be plausible alike, a third of the automatic choices
differ from the user, but EX drops by only 0.4 points, so most of these
choices do not change whether the query is correct.

\textbf{Per type.} Table~\ref{tab:pertype} breaks the automatic
resolutions at $\theta^{*}$ down by type. The scorer leaves the types
that change the row set to the user. Of the 156 Join, NULL, and Tie
ambiguities on the three benchmarks, it auto-resolves 7, all in
agreement with the simulated user, and it auto-resolves at most 11\%
of the Scope ambiguities on any benchmark. Most automatic resolutions
are Column, Project, and Aggregate ambiguities. The low agreement on
CLAMBSQL comes mostly from Column ambiguities, which account for 795
of its 1{,}339 automatic resolutions and agree with the user on
62.9\%. CLAMBSQL constructs its column
ambiguities so that either column fits the question, and the scorer,
trained on BIRD databases only, has little to tell them apart.

\begin{table}[t]
\centering
\caption{Automatic resolutions per type at $\theta^{*}$ (\coder). Auto: auto-resolved out of detected. Agr.: agreement with the simulated user (\%).}
\label{tab:pertype}
\footnotesize
\setlength{\tabcolsep}{3pt}
\begin{tabular}{@{}lrrrrrr@{}}
\toprule
& \multicolumn{2}{c}{BIRD} & \multicolumn{2}{c}{AMBROSIA} &
\multicolumn{2}{c}{CLAMBSQL} \\
\cmidrule(lr){2-3}\cmidrule(lr){4-5}\cmidrule(lr){6-7}
Type & Auto & Agr. & Auto & Agr. & Auto & Agr. \\
\midrule
Scope & 19/172 & 78.9 & 13/859 & 84.6 & 11/613 & 72.7 \\
Attachment & 1/49 & 100 & 6/144 & 66.7 & 6/113 & 100 \\
Column & 249/487 & 92.0 & 486/1{,}271 & 93.4 & 795/1{,}822 & 62.9 \\
Join & 1/16 & 100 & 3/46 & 100 & 2/37 & 100 \\
Value & 27/144 & 77.8 & 25/179 & 96.0 & 37/239 & 75.7 \\
NULL & 0/4 & -- & 0/1 & -- & 0/0 & -- \\
Aggregate & 96/352 & 92.7 & 43/186 & 100 & 135/690 & 71.9 \\
Project & 142/526 & 94.4 & 292/1{,}279 & 95.5 & 353/1{,}489 & 73.7 \\
Tie & 0/22 & -- & 1/10 & 100 & 0/20 & -- \\
\bottomrule
\end{tabular}
\end{table}

\textbf{Ablation.}\label{sec:exp-rq3} Table~\ref{tab:ablation} adds one component at a
time on BIRD, whose real databases exercise all checks.
 Without the
ambiguity checks, forward detection asks \abAblQFwdNoCheck{} questions
per instance and reaches 61.8\% EX. The checks remove 31\% of these
questions and add 3.2 points, because the
existence and difference checks write values verified on $D$ that
the user's answers do not supply. Backward detection adds 1.4
points at 0.42 more questions per instance, and 2.2 and 2.0 points on
AMBROSIA and CLAMBSQL. The scorer then removes 35\% of the
questions at unchanged EX. Auto-resolving as many ambiguities at random
lowers EX to 64.9\%, so the scorer's choice
preserves accuracy.

% \begin{table}[t]
% \centering
% \caption{Ablation on BIRD (\coder). EX under $y^{*}$ (\%). Each row
% adds one component to the row above, and arrows give the change over
% that row. Rows above the scorer ask every remaining ambiguity.
% $^{\dagger}$Replaces the scorer: auto-resolves as many ambiguities as
% the scorer, chosen at random with a random interpretation (mean of
% three seeds), arrows relative to the backward-detection row.}
% \label{tab:ablation}
% \footnotesize
% \begin{tabular}{@{}lcc@{}}
% \toprule
% Variant & EX & $Q$ \\
% \midrule
% Never-Ask & \abBirdZero & 0 \\
% Forward detection, no checks & \abAblFwdNoCheck & \abAblQFwdNoCheck \\
% \quad + ambiguity checks & \abAblFwd & \abAblQFwd \\
% \quad + backward detection & \abAblAskAll & \abAblQAskAll \\
% \quad + scorer & \abAblJevAsk & \abAblQJevAsk \\
% \quad + random auto-resolution$^{\dagger}$ & \abAblRand & \abAblQRand \\
% \bottomrule
% \end{tabular}
% \end{table}


\begin{table}[t]
\centering
\caption{Ablation on BIRD (\coder). Each row adds one component, and arrows give the change over the row above.}
\label{tab:ablation}
\footnotesize
\begin{tabular}{@{}lcc@{}}
\toprule
Variant & EX & $Q$ \\
\midrule
Never-Ask & \abBirdZero & 0 \\
Forward detection & \abAblFwdNoCheck & \abAblQFwdNoCheck \\
\quad + ambiguity checks & \abAblFwd & \abAblQFwd \\
\quad + backward detection & \abAblAskAll & \abAblQAskAll \\
\quad + scorer & \abAblJevAsk & \abAblQJevAsk \\
\bottomrule
\end{tabular}
\end{table}


%======================================================================
\subsection{User Study (RQ3)}
\label{sec:exp-user}

The results above rely on a simulated user. To test whether real users
answer our questions as usefully, we replace it with people on a
subset of the benchmarks, with \terra{} as the parser.

\textbf{Protocol.} We sample 100 questions, 40 from BIRD and 30 each
from AMBROSIA and CLAMBSQL, among those where Ours or AmbiSQL asks at
least one question, spread round-robin over the databases with a fixed
seed. \pending{Ten} computer science graduate students with at least
one year of SQL experience take part. None is an author or familiar
with the benchmarks. Each answers 10 questions, and each question is
handled by exactly one participant. A participant receives the same
information as the simulated user: the question, the schema with
column descriptions and example values, and an intent card with the
gold query and a one-sentence description of the intended result. The
two systems are shown anonymously, in counterbalanced order. For every
clarification question, the participant selects an option or answers
``Others'' in free text, and rates whether it is clear and
whether it helps to obtain the intended result, both on a 1--5 scale.
For the ambiguities that our scorer resolves automatically, the
participant states whether the written interpretation is the one they
mean. We write the answers into the evidence with the same procedure
as for the simulated user and reparse with \terra.

\begin{table}[t]
\centering
\caption{User study on 100 questions (\terra). EX (\%) with simulated (Sim.) and participants' (Part.) answers and McNemar $p$. \cmark\xmark: correct with Sim.\ only, \xmark\cmark: with Part.\ only. Agr.: same choice as the simulated user (\%).}
\label{tab:user}
\footnotesize
\setlength{\tabcolsep}{3pt}
\begin{tabular}{@{}lccccccc c@{}}
\toprule
Method & Sim. & Part. & $p$ & \cmark\cmark & \cmark\xmark & \xmark\cmark &
\xmark\xmark & Agr. \\
\midrule
Never-Ask & 60.0 & -- & -- & -- & -- & -- & -- & -- \\
AmbiSQL~\cite{ding2025ambisql} & 70.0 & \pending{--} & \pending{--} & \pending{--} &
\pending{--} & \pending{--} & \pending{--} & \pending{--} \\
Ours & 78.0 & \pending{--} & \pending{--} & \pending{--} &
\pending{--} & \pending{--} & \pending{--} & \pending{--} \\
\bottomrule
\end{tabular}
\end{table}

\textbf{Results.} With the simulated user, Ours reaches 78.0\% EX on
the sample, against 70.0\% for AmbiSQL and 60.0\% for Never-Ask. Reparsing
the same answers once more gives 75.0\%, 69.0\%, and 62.0\%, so the
ranking holds although \terra{} is not deterministic.
\pending{With the participants' answers, Ours reaches X\% against Y\%
for AmbiSQL ($p=\cdot$). Relative to the simulated user, EX changes by
Z points for Ours and by W points for AmbiSQL, neither significant
($p=\cdot$ and $\cdot$). The off-diagonal cells of
Table~\ref{tab:user} are small and roughly symmetric, so the simulated
user neither systematically overestimates nor underestimates the
accuracy that real users reach. The participants choose the same
option as the simulated user on A\% of our questions and B\% of
AmbiSQL's, and accept C\% of our automatic resolutions. They rate our
questions D and AmbiSQL's E for clarity, and F and G for usefulness.}

%======================================================================
\subsection{Case Study}
\label{sec:exp-case}

\textbf{Backward detection.} The BIRD question ``Among the customers
who paid in euro, how many of them have a monthly consumption of over
1000?'' gets forward questions on how monthly consumption is computed,
but none on what is counted, and the parser counts monthly records
with \texttt{COUNT(*)}. The drafted queries disagree on this point,
and backward detection asks whether to count distinct customers or
monthly records. With the answer, the parser writes
\texttt{COUNT(DISTINCT CustomerID)}, as the gold query does.

\textbf{Scope.} For the AMBROSIA question ``List the animals we have
in each zoo,'' AmbiSQL asks which SQL operation to use. All its
options (count per zoo, list per zoo, distinct animal types per zoo)
read ``each zoo'' distributively, so the simulated user picks the
closest ones, and AmbiSQL returns the query of Never-Ask. Our Scope
ambiguity also offers the collective reading, the animals present in
every zoo. With it, the parser groups by animal and keeps those whose
count of distinct zoos equals the number of zoos, as the gold query
does.

\textbf{A scorer error.} The question ``show me all
projects by company'' occurs in AMBROSIA and in CLAMBSQL with the same
two readings of \emph{company}, the real estate developer and the
contractor. The gold query uses the developer in AMBROSIA and the
contractor in CLAMBSQL. The scorer resolves both to the developer
($H_{\mathrm{norm}}=0.42$ and $0.39$), which is wrong on CLAMBSQL.

%======================================================================
\subsection{Limitations}
\label{sec:discussion}


\textbf{Failures of clarification.}\label{sec:exp-fail} Asking rarely
hurts. Among the questions where the simulated user was asked and the
written interpretation matches the answer, asking turns a correct
query into a wrong one on at most 1.6\% on AMBROSIA and 3.0\% on
CLAMBSQL, against 394 and 184 questions that it turns from wrong to
correct. The remaining errors have three sources. First, the options
may not cover the reading of the gold query. On one CLAMBSQL question,
both options name the car maker while the gold query returns its
identifier. Second, the parser may ignore a correct clarification. For
the BIRD question ``Name the user that commented `thank you
user93!','' the evidence names the display name of the user, but every
system reads the near-identical column of the comments table, which is
empty for this row. Third, an automatic resolution of two equally
plausible readings can only follow a prior
(Section~\ref{sec:exp-case}).

\textbf{Soundness of not asking.} The queries that the difference
check executes for Scope and Attachment are written by an LLM and may
not implement their readings, although a second call verifies them
(Section~\ref{sec:method-check}). The relevance check tests one
ambiguity at a time and can miss ambiguities that change the result
only jointly, for example when the chosen column changes the values of
another ambiguity. It also probes \coder{} for every parser, so for
DeepEye-SQL and \terra{} it predicts relevance from the behavior of
\coder{} (Section~\ref{sec:exp-setup}). A missed ambiguity costs
accuracy only when the parser then picks the wrong reading, and
Table~\ref{tab:terra} shows that the gains persist for both parsers.

\textbf{Generalization.} The scorer is trained and evaluated with a
simulated user who knows the gold query, as in prior
work~\cite{zhang2025clear,xu2026visinteract}. A user who holds only
domain knowledge may answer less reliably or reject every option, and
our user study (Section~\ref{sec:exp-user}) covers 100 questions.
The scorer is trained on BIRD databases only. On CLAMBSQL, whose
column ambiguities are constructed so that either column fits, its
agreement with the simulated user falls to 62.9\% for Column
(Table~\ref{tab:pertype}). A threshold selected on validation data
meets the target agreement only when the test ambiguities resemble
the validation ones, an assumption shared by selective
prediction~\cite{geifman2017selective,kamath2020selective}. Finally,
AMBROSIA and CLAMBSQL contain few data-dependent ambiguities, which
the existence and difference checks target, so the checks are
exercised mainly on BIRD (Table~\ref{tab:ablation}).

\textbf{Practical cost.} Fewer questions to the user come at the price
of more computation. Forward detection makes three LLM calls, backward
detection drafts six queries and makes a final call
(Section~\ref{sec:method-detect}), and the relevance check makes
\probeCallsBird, \probeCallsAmbrosia, and \probeCallsClamb{} \coder{}
calls per question on BIRD, AMBROSIA, and CLAMBSQL. The scorer adds
one forward pass of Qwen3-8B per ambiguity and requires no decoding.
These calls add latency before the first question reaches the user.
All checks also need read-only access to the database, as do other
methods that probe the database~\cite{tian2026pv,xiong2025multi}.
In exchange, Ours is more accurate than CLEAR and AmbiSQL on all
three benchmarks while asking about as many questions as AmbiSQL
(Table~\ref{tab:main}).
